\documentclass[twocolumn]{iopjournal}

\usepackage{amsmath,amssymb,amsfonts}

\fancyhfoffset[L]{0mm}
\setlength{\marginparwidth}{0pt}
\setlength{\marginparsep}{0pt}

% ==========================================
% ZENODO PREPRINT HEADERS (ACTIVE)
% ==========================================
\renewcommand{\articletype}[1]{%
  {\vspace*{-8mm}\noindent \Large \sf Preprint}\par
  \vspace*{6mm}\noindent{\scriptsize \sf{\bfseries \MakeUppercase{#1}}}}

\fancyhead[L]{}
\fancyhead[R]{}

\hypersetup{pdftitle={Supplementary Material: The Constitutive Law from Exchange Symmetry}, pdfauthor={Stephen Nelson}, pdfsubject={Gravity Field Dynamics (GFD)}, pdfkeywords={constitutive law, interpolating function, MOND, AQUAL, exchange symmetry, algebra, proof}}

\begin{document}

\articletype{Supplementary Material}

\title{Supplementary Material: The Constitutive Law\\
from Exchange Symmetry}

\author{Stephen Nelson}

\affil{Field Dynamics Research}

\ead{stephen@fielddynamics.org}

\begin{abstract}
The tetrad equation produces two parity-related bounded regions,
TetA and TetB~\cite{Nelson2026geometry}. Let the acceleration and
the fraction $\mu$ of field capacity carried by TetA satisfy a
polynomial equation. Three premises determine that equation: the endpoint $\mu(0) = 0$; closure of the partition states
under the exchange of the regions; and degree one in each variable,
the degree in $\mu$ counting the directions in which the coupling
runs. A short lemma on degrees makes the relation bilinear, and the
premises then force the screening identity $x(1-\mu) = \mu$, that is
$\mu(x) = x/(x+1)$. The exchange is derived, not assumed: it is the
inversion $x \to 1/x$ about the scale $a_0$, and parity reversing the
coordinate $s = 2\mu - 1$. Removing any premise admits explicit
alternatives, among them the framework's own bound-state law, so each
is necessary. No coupling polynomial, field equation or equilibrium
postulate enters the proof.
\end{abstract}

\keywords{constitutive law, interpolating function, MOND, AQUAL, exchange symmetry, algebra, proof}

% ----------------------------------------------------------------------

\section*{Definitions}

Let $\tilde x = g/g_\ast \in [0,\infty]$ be the acceleration in units
of an arbitrary reference $g_\ast$; the natural unit is fixed below.
Let $\mu(\tilde x) \in [0,1]$ be the fraction of the total field
capacity carried by TetA, so that $1-\mu$ is the fraction carried by
TetB, and let
\begin{equation}
  \Omega = \frac{\mu}{1-\mu} \in [0,\infty]
  \label{eq:odds}
\end{equation}
be the odds ratio between the two regions.

% ----------------------------------------------------------------------

\section*{Premises}

\paragraph{(E) Endpoint.} $\mu(0) = 0$: at zero acceleration the
field sits entirely in TetB.

\paragraph{(X) Exchange closure.} The main
letter~\cite{Nelson2026geometry} and its
supplement~\cite{Nelson2026scheme} show that the identity and parity
components of the tetrad equation give TetA and TetB, exchanged by
spatial parity $P$. Parity leaves $|g|$ unchanged; it acts on the
partition states, sending the state in which TetA carries a share
$\mu$ to the state in which TetB carries that share. Since $P$ is a
symmetry, the set of partition states realised on $[0,\infty]$ is
closed under it: if $\mu$ is realised at some acceleration, so is
$1 - \mu$. Nothing more about the exchange is assumed.

\paragraph{(A) Degree one.} The tetrad equation is polynomial, and
we take $\tilde x$ and $\mu$ to satisfy a real polynomial equation
$P(\tilde x, \mu) = 0$, $P$ irreducible. We
require $P$ to have degree one in $\mu$ and degree one in
$\tilde x$. The degree in $\mu$ counts the directions in which the
coupling runs. The unbound field couples in one direction: matter
sources the field, and the field does not source
itself~\cite{Nelson2026bound}. A bound field sources itself, couples
in both directions, and obeys $x(1-\mu)^2 = \mu$, of degree two in
$\mu$~\cite{Nelson2026bound}. The degree in $\tilde x$ is the
exponent of the driving term of the second
route~\cite{Nelson2026equilibrium}.

% ----------------------------------------------------------------------

\section*{Lemma}

\paragraph{Lemma.} \textit{Let $K$ be a field and $P \in K[x,y]$
irreducible, of degree $m \ge 1$ in $y$ and $n \ge 1$ in $x$, with
$P(x,y) = 0$. Then $[K(x,y):K(x)] = m$ and $[K(x,y):K(y)] = n$. So
$y \in K(x)$ and $x \in K(y)$ if and only if $m = n = 1$, that is,
$P = axy + bx + cy + d$ with $ad - bc \neq 0$.}

\paragraph{Proof.} $P$ is irreducible in $K[x][y]$ and has positive
degree in $y$, so it is primitive over $K[x]$, and by Gauss's lemma
it is irreducible in $K(x)[y]$. It is therefore the minimal
polynomial of $y$ over $K(x)$, of degree $m$; by symmetry
$[K(x,y):K(y)] = n$. Finally, $P = (ay + b)x + (cy + d)$ factors
exactly when $(a,b)$ and $(c,d)$ are proportional, that is when
$ad - bc = 0$.
\hfill $\square$

% ----------------------------------------------------------------------

\section*{The Unique Law}

\paragraph{Theorem.} \textit{Under (E), (X) and (A), with
$x = g/a_0$ for a unique scale $a_0$, the partition relation is
$x(1-\mu) = \mu$, so that $\mu(x) = x/(x+1)$, and the exchange acts
on accelerations as $x \to 1/x$.}

\paragraph{Step 1: bilinear.} By (A) and the lemma,
$P = a\tilde x\mu + b\tilde x + c\mu + d$ with $ad - bc \neq 0$. By
(E), $P(0,0) = d = 0$, so $bc \neq 0$, and in odds
\begin{equation}
  \Omega(\tilde x) = \frac{\alpha \tilde x}{\gamma \tilde x + \delta},
\end{equation}
with $\alpha = -b$, $\gamma = a + b$, $\delta = c$ and
$\alpha\delta \neq 0$.

\paragraph{Step 2: the pole.} By (E) the state $\mu = 0$ is realised,
so by (X) the state $\mu = 1$, $\Omega = \infty$, is realised at
some $\tilde x \in [0,\infty]$. The only pole of $\Omega$ is
$\tilde x = -\delta/\gamma$, read as $\infty$ when $\gamma = 0$. If
$\gamma\delta < 0$ the pole is positive and simple, and $\Omega$
changes sign across it, contrary to $\Omega \ge 0$. If
$\gamma\delta > 0$ there is no pole on $[0,\infty]$, since
$\Omega(\infty) = \alpha/\gamma$ is finite, contrary to (X). Hence
$\gamma = 0$, $a = -b$, and $\Omega = k\tilde x$ with
$k = \alpha/\delta > 0$ because $\Omega \ge 0$.

\paragraph{Step 3: the scale.} Setting $a_0 = g_\ast/k$ and
$x = g/a_0$ gives $\Omega(x) = x$, and by Eq.~\eqref{eq:odds}
\begin{equation}
  \boxed{x\,(1-\mu) = \mu, \qquad \mu(x) = \frac{x}{x+1}.}
  \label{eq:law}
\end{equation}
The map $(x,\mu) \to (1/x,\, 1-\mu)$ carries Eq.~\eqref{eq:law} to
itself, and $\mu$ is single-valued in $x$, so the state exchanged
with the one at $x$ sits at $1/x$: the exchange is the inversion
$\iota(x) = 1/x$, whose fixed point is $x = 1$, that is $g = a_0$.
\hfill $\square$

\paragraph{Consequences.} The argument identifies $a_0$ as the fixed
point of the exchange without evaluating it. The second endpoint,
$\mu \to 1$ as $x \to \infty$, is derived. The complement symmetry
$\mu(x) + \mu(1/x) = 1$ is derived. The involution property
$\iota \circ \iota = \mathrm{id}$, which parity requires through
$P^2 = 1$, is not used and is recovered. At the fixed point
$\mu(1) = 1/2$.

\paragraph{Parity form.} Let $s = 2\mu - 1 \in [-1,1]$, so that TetA
is $s = +1$ and TetB is $s = -1$. Equation~\eqref{eq:law} becomes
\begin{equation}
  x\,(1-s) = 1+s, \qquad s(x) = \frac{x-1}{x+1},
\end{equation}
and the exchange becomes $(x, s) \to (1/x, -s)$. In this coordinate
the exchange is parity itself, reversing $s$ as it reverses the
rank-one term $Ks$ of the coupling polynomial
$f(K,s) = 1 + sK + K^2$, whose normalisation
$\big(f(K,s) - f(K,-1)\big)/2K$ is
$\mu = (1+s)/2$~\cite{Nelson2026geometry}.

% ----------------------------------------------------------------------

\section*{Necessity of the Premises}

Removing any one premise admits alternatives satisfying the other
two:

\vspace{6pt}
\noindent
\begin{tabular*}{\columnwidth}{@{\extracolsep{\fill}}lp{0.74\columnwidth}@{}}
\hline
Removed & Admitted alternatives \\
\hline
(E) & $\Omega = (x+b)/(bx+1)$, $0 < b < 1$; here $\mu(0) = b/(1+b)$
  and the states fill $[\,b/(1+b),\,1/(1+b)\,]$ \\
(X) & $\Omega = \alpha x/(\gamma x + \delta)$, $\alpha, \gamma,
  \delta > 0$; here $\mu(\infty) < 1$ \\
(A) & $x^n(1-\mu)^m = \mu$, $n, m \ge 1$, $(n,m) \neq (1,1)$, of
  degree $n$ in $x$ and $m$ in $\mu$; $(n,m) = (1,2)$ is the
  bound-state law~\cite{Nelson2026bound} \\
\hline
\end{tabular*}
\vspace{6pt}

Each relation in the (A) row is irreducible (Eisenstein at the prime
$\mu$ of $\mathbb{R}[\mu]$), has a single root
$\mu \in [0,1]$ for each $x$, and runs from $\mu = 0$ to $\mu = 1$.
Without (A) the exchange carries no content: any continuous,
strictly increasing $\mu$ with $\mu(0) = 0$ and $\mu \to 1$
satisfies (E) and (X).

% ----------------------------------------------------------------------

\section*{Relation to the Other Routes}

The main letter derives the same law by two further routes: graded
rank counting of the tetrad equation, with parity interpolation and
rational normalisation~\cite{Nelson2026metric}; and the AQUAL field
equation with single-channel detailed
balance~\cite{Nelson2026equilibrium}. The proof above uses neither.
The three routes meet in
one relation of degree one in each variable. In the first it is the
map $s(x)$: of its three boundary conditions, $s(0) = -1$ is (E),
$s(1) = 0$ is the choice of unit and $s(\infty) = +1$ is derived,
and its fractional linear form is (A) through the lemma. In the
second it is the balance $x^n(1-\epsilon) = \epsilon/K$, of degree
$n$ in $x$ and one in $\epsilon$, with $n = 1$ fixed there by
deep-field scaling.

% ----------------------------------------------------------------------

\section*{Summary}

For a coupling of degree one in each variable, closure of the
partition states under the exchange of TetA and TetB forces the
screening identity $x(1-\mu) = \mu$, whose solution is the
constitutive law $\mu(x) = x/(x+1)$; each premise is necessary.

This function was identified empirically by
Milgrom~\cite{Milgrom1983} and selected by fitting by
Famaey and Binney~\cite{FamaeyBinney2005}. The present
derivation obtains it from the exchange symmetry of the
two parity-related roots of the equation that defines the
gravitational field.

% ----------------------------------------------------------------------

\begin{thebibliography}{99}

\bibitem{Nelson2026geometry}
S.~Nelson,
\textit{The Coordinate-Space Geometry of the Tetrad Frame},
Class.\ Quantum Grav.\ (2026), submitted.

\bibitem{Nelson2026scheme}
S.~Nelson,
\textit{Supplementary Material: The Tetrad Scheme and the Two
Bounded Regions} (2026).

\bibitem{Nelson2026bound}
S.~Nelson,
\textit{Nuclear Binding Energy from the Bidirectional
Constitutive Law} (2026).

\bibitem{Nelson2026metric}
S.~Nelson,
\textit{Supplementary Material: The Constitutive Law Read from
the Metric} (2026).

\bibitem{Nelson2026equilibrium}
S.~Nelson,
\textit{Supplementary Material: The Constitutive Law from
Equilibrium Field Closure} (2026).

\bibitem{Milgrom1983}
M.~Milgrom,
\textit{A modification of the Newtonian dynamics as a possible
alternative to the hidden mass hypothesis},
Astrophys.\ J.\ \textbf{270} (1983) 365.

\bibitem{FamaeyBinney2005}
B.~Famaey and J.~Binney,
\textit{Modified Newtonian dynamics in the Milky Way},
Mon.\ Not.\ R.\ Astron.\ Soc.\ \textbf{363} (2005) 603.

\end{thebibliography}

\end{document}
